% !TeX root = ../master.tex

\lecture{7}{2026-09-15}{Makromekanik af Lamina, Pt. II}

\subsection{Arbitrær orientering}
For fiberretninger og belastninger med ``arbitrære'' lokationer må vi transformere fra det lokale $(1,2,3)$-koordinatsystem til et globalt $(x,y,z)$ koordinatsystem:
\begin{equation*}
  \begin{bmatrix}
    \sigma_x\\
  \sigma_y\\
  \tau_{xy}\\
\end{bmatrix} = [T]^{-1} \begin{bmatrix}
  \sigma_1\\
\sigma_2\\
\tau_{12}\\
\end{bmatrix}
\end{equation*}
hvor $\left[ T \right]$ er tranformationsmatricen bestående af $\cos$ og $\sin$ led samt vinklen $\theta$ mellem $(1,2,3)$- og $(x,y,z)$-systemerne.

Det samme kan gøres for tøjninger idet vi bruger $\epsilon_{12}$
\begin{equation*}
  \begin{bmatrix}
    \epsilon_x\\
  \epsilon_y\\
  \frac{1}{2}\gamma_{xy}\\
  \end{bmatrix} = \left[ T \right]^{-1} \begin{bmatrix}
    \epsilon_1\\
  \epsilon_2\\
  \frac{1}{2}\gamma_{12}\\
  \end{bmatrix}
.\end{equation*}
Idet vi introducerer Reuter matricen $R$:
\begin{equation*}
  \left[ R \right] = \begin{bmatrix}
    1 & 0 & 0\\
  0 & 1 & 0\\
  0 & 0 & 2\\
  \end{bmatrix}
\end{equation*}
kan $\frac{1}{2}$ fjernes:
\begin{equation*}
  \begin{bmatrix}
    \epsilon_x\\
  \epsilon_y\\
  \gamma_{xy}\\
  \end{bmatrix} = \left[ R \right]\left[ T \right]^{-1} \left[ R \right]^{-1} \begin{bmatrix}
    \epsilon_1\\
  \epsilon_2\\
  \gamma_{12}\\
  \end{bmatrix}
.\end{equation*}

Spændings-tøjningerelationen for et lamina med $\theta$ rotation er
\begin{equation*}
  \begin{bmatrix}
    \sigma_x\\
  \sigma_y\\
  \tau_{xy}\\
\end{bmatrix} = \left[ T \right]^{-1} \left[ Q \right] \underbracket{\left[ R \right]\left[ T \right]\left[ R \right]^{-1}}_{\left[ T \right]^{-T}} \begin{bmatrix}
  \epsilon_x\\
\epsilon_y\\
\gamma_{xy}\\
\end{bmatrix}
\end{equation*}
eller
\begin{align*}
  \begin{bmatrix}
    \sigma_x\\
  \sigma_y\\
  \tau_{xy}\\
  \end{bmatrix} &= \left[ T \right]^{-1} \left[ Q \right] \underbracket{\left[ T \right]^{-1} \left[ Q \right]\left[ T \right]^{-T}}_{\left[ \overline{Q} \right]} \begin{bmatrix}
  \epsilon_x\\
\epsilon_y\\
\gamma_{xy}\\
\end{bmatrix} \\
\begin{bmatrix}
  \sigma_x\\
\sigma_y\\
\tau_{xy}\\
\end{bmatrix} &= \begin{bmatrix}
  \overline{Q}_{11} & \overline{Q}_{12} & \overline{Q}_{16}\\
\overline{Q}_{12} & \overline{Q}_{22} & \overline{Q}_{26}\\
\overline{Q}_{16} & \overline{Q}_{26} & \overline{Q}_{66}\\
\end{bmatrix} \begin{bmatrix}
  \epsilon_x\\
\epsilon_y\\
\gamma_{xy}\\
\end{bmatrix}
\end{align*}
hvor $[\overline{Q}]$ er den transformerede reducerede planstress stivhedsmatrice.

